% Figure: temperature against altitude for a rising air parcel. The dry
% adiabatic lapse rate (about 9.8 K/km) is the steeper line followed by an
% unsaturated parcel; the parcel's dew-point temperature falls much more
% slowly with height (about 1.8 K/km). The two lines converge at the lifting
% condensation level, the cloud base, where the parcel reaches saturation
% and a cloud forms. Above the cloud base a saturated parcel follows the
% gentler moist adiabatic lapse rate (about 6 K/km) because condensation
% releases latent heat. Coordinates are illustrative of the standard
% construction for a parcel starting at 25 C surface with a 12 C dew point.
\begin{figure}[htbp]
\centering
\begin{tikzpicture}
\begin{axis}[
  width=11cm, height=8.5cm,
  xlabel={temperature ($^{\circ}\mathrm{C}$)},
  ylabel={altitude $z$ (km)},
  xmin=-10, xmax=28, ymin=0, ymax=4,
  axis lines=left,
  tick label style={font=\scriptsize},
  label style={font=\small},
  legend pos=north east,
  legend style={font=\scriptsize},
]
% Dry adiabat: parcel temperature 25 C at surface, slope -9.8 K/km.
% At cloud base z = 1.625 km, T = 25 - 9.8*1.625 = 9.1 C.
\addplot[engred, very thick] coordinates {
  (25,0) (9.075,1.625)
};
\addlegendentry{dry adiabat ($9.8\,\text{K/km}$)}
% Dew-point line: 12 C at surface, slope -1.8 K/km, meets parcel at z=1.625.
% At z=1.625, dewpoint = 12 - 1.8*1.625 = 9.075 C. Lines meet here.
\addplot[engblue, very thick, dashed] coordinates {
  (12,0) (9.075,1.625)
};
\addlegendentry{dew point ($1.8\,\text{K/km}$)}
% Moist adiabat above the cloud base: slope -6 K/km from (9.075,1.625).
% At z = 4 km: T = 9.075 - 6*(4 - 1.625) = 9.075 - 14.25 = -5.175 C.
\addplot[engamber, very thick] coordinates {
  (9.075,1.625) (-5.175,4)
};
\addlegendentry{moist adiabat ($6\,\text{K/km}$)}
% mark the lifting condensation level (cloud base)
\addplot[only marks, mark=*, engblack, mark size=2.2pt] coordinates {(9.075,1.625)};
\node[engblack, font=\scriptsize, anchor=west] at (axis cs:9.5,1.625) {cloud base (LCL)};
\end{axis}
\end{tikzpicture}
\caption[Adiabatic lapse rate and cloud base]{Temperature against altitude
for an air parcel lifted from a $25\,\degC$ surface with a $12\,\degC$ dew
point. The unsaturated parcel cools along the dry adiabatic lapse rate of
about $9.8\,\text{K/km}$ (red), while its dew-point temperature, the
saturation temperature of the water it carries, falls far more slowly at
about $1.8\,\text{K/km}$ (blue dashed) because lifting barely changes the
vapour content. The two lines meet at the lifting condensation level near
$1.6\,\text{km}$, the cloud base, where the parcel reaches saturation and
its water vapour begins to condense. Above the cloud base the parcel
follows the gentler moist adiabatic lapse rate of about $6\,\text{K/km}$
(amber) because the latent heat released by condensation offsets part of
the expansion cooling. The construction is the equation of state and the
saturation curve acting together: the ideal-gas law sets the dry cooling
rate, and the Clausius-Clapeyron saturation curve sets the height at which
condensation begins.}
\label{fig:vol04ch02:lapse-rate-cloud-base}
\end{figure}
